%=======================================================================
%  CHAPTER 3 — RADIO WAVE PROPAGATION IN THE MOBILE ENVIRONMENT  (12 hrs)
%=======================================================================
\section{Radio Wave Propagation in Mobile Network Environment}

{\color{sub}\footnotesize 12 hrs, \mk{the largest chapter in the syllabus}
\gap$\cdot$\gap asked in \textbf{22 of 22 papers} \gap$\cdot$\gap 7 syllabus sub-topics
\gap$\cdot$\gap \mk{by far the heaviest numerical load} $\rightarrow$ Ch 3 numerical companion.}

\lead{Read this first:}
\begin{itemize}
  \item The chapter splits cleanly in two, and the exam treats them separately.
  \begin{itemize}
    \item \textbf{Large-scale path loss}: how far the signal reaches, averaged over
          hundreds of wavelengths. Free space, two-ray, Okumura, Hata, indoor.
    \item \textbf{Small-scale fading}: how the signal flutters over a few wavelengths.
          Delay spread, Doppler, coherence bandwidth, Rayleigh / Ricean.
  \end{itemize}
  \item Three questions are near-certain:
        \textbf{small-scale fading types + factors} \tS{6},
        \textbf{coherence bandwidth and coherence time} \tF{5}, and an
        \textbf{Okumura or Hata numerical} \tF{4}.
  \item Two derivations are worth memorising line by line, because they recur:
        the \textbf{two-ray path difference and phase difference} \tF{4} and the
        \textbf{Doppler shift} \tF{3}.
  \item \mk{Nearly every question here has a numerical twin.} Learn the theory with the
        companion open.
\end{itemize}

\hr

\T{3.1 Free space propagation}

\Q{Explain free space propagation. Relate radiated power to electric field}{\m{8}\m{2+2+2+2}\m{6}}
{\color{sub}\footnotesize \tF{3} \yr{\texttt{80 Ba}, \texttt{81 Ba}, \textbf{78 Ch},
\textbf{\texttt{80 Bh}}} \gap$\cdot$\gap \mk{almost always set as a numerical}
$\rightarrow$ Ch 3 companion}

\sbs{%
\lead{When it applies}
\begin{itemize}
  \item Used to predict received signal strength when Tx and Rx have a
        \textbf{clear, unobstructed LOS path}.
  \item Eg satellite links, microwave LOS links.
  \item \mk{Almost never true for a mobile in a city}, but every other model is built as a
        correction to it.
\end{itemize}

\lead{Step 1: power flux density}
\begin{itemize}
  \item An \textbf{isotropic} source radiates $P_t$ uniformly over a sphere of area $4\pi d^2$.
  \item \textbf{Power flux density} at distance $d$:
        \[ P_d \;=\; \frac{P_t G_t}{4\pi d^2} \quad \text{W/m}^2 \]
  \item \textbf{EIRP} $= P_t G_t$ watts: the power an isotropic antenna would need to give the
        same reading in the main beam.
  \item \textbf{ERP} uses a half-wave dipole as the reference instead:
        $\text{ERP} = \text{EIRP} - 2.15$ dB.
\end{itemize}}{%
\lead{Step 2: effective aperture}
\begin{itemize}
  \item \textbf{Effective aperture} $A_e$ $=$ ratio of power available at the antenna terminals to
        the incident power per unit area.
        \[ A_e \;=\; \frac{G_r \lambda^2}{4\pi}, \qquad \lambda = \frac{c}{f} \]
\end{itemize}

\lead{Step 3: the Friis equation}
\begin{itemize}
  \item $P_r = P_d \cdot A_e$, so
        \[ \boxed{\;P_r(d) \;=\; \frac{P_t G_t G_r \lambda^2}{(4\pi)^2 d^2 L}\;} \]
  \item $L \ge 1$ $=$ system loss (feeder attenuation, filters, mismatch). $L=1$ means
        no hardware loss.
  \item \mk{$P_r \propto 1/d^2$ and $\propto \lambda^2$}: doubling distance costs 6 dB,
        doubling frequency costs 6 dB.
\end{itemize}}

\lead{Path loss}
\begin{itemize}
  \item \textbf{With} antenna gains:
        \[ PL(\text{dB}) = 10\log\frac{P_t}{P_r}
           = -10\log\!\left[\frac{G_t G_r \lambda^2}{(4\pi)^2 d^2}\right] \]
  \item \textbf{Without} antenna gains ($G_t = G_r = 1$):
        \[ PL(\text{dB}) = -10\log\!\left[\frac{\lambda^2}{(4\pi)^2 d^2}\right] \]
  \item Useful dB form, $d$ in km and $f$ in MHz:
        \[ PL_{\text{FS}} = 32.44 + 20\log f_{\text{MHz}} + 20\log d_{\text{km}} \]
        \added
  \item The minus sign on the gain term says that term is a \textbf{gain}. The
        $(4\pi d/\lambda)^2$ term is the \textbf{free space loss} proper.
  \item Raising $d$ raises the loss, which pushes designers to \textbf{lower frequencies} to keep
        the link budget manageable.
\end{itemize}

\penalty0\vspace{2pt}
\sbs{%
\lead{Relating power to E-field \mk{(the 78 Ch / 80 Ba style numerical)}}
\begin{itemize}
  \item Power flux density in terms of the field:
        \[ P_d \;=\; \frac{|E|^2}{\eta} \;=\; \frac{|E|^2}{120\pi}, \qquad
           \eta = 377\;\Omega \text{ (free space)} \]
  \item So $E$ falls as $1/d$ while power falls as $1/d^2$.
  \item Received power through the aperture:
        \[ P_r(d) \;=\; \frac{|E|^2}{120\pi}\cdot \frac{G_r\lambda^2}{4\pi} \]
  \item \textbf{rms voltage} into a matched receiver of input resistance $R_{\text{ant}}$:
        \[ P_r = \frac{V_{\text{rms}}^2}{R_{\text{ant}}} = \frac{(V_{\text{ant}}/2)^2}{R_{\text{ant}}}
           \;\Longrightarrow\; V_{\text{ant}} = 2\sqrt{P_r R_{\text{ant}}} \]
        \mk{The factor 2 is the matched-load voltage divider. Losing it is the standard mistake.}
\end{itemize}}{%
\lead{Far field (Fraunhofer region)}
\begin{itemize}
  \item Friis is valid only \textbf{beyond} the far-field distance
        \[ d_f \;=\; \frac{2D^2}{\lambda} \]
        where $D$ $=$ the largest physical linear dimension of the antenna.
  \item Also needs $d_f \gg D$ and $d_f \gg \lambda$.
  \item The equation blows up at $d = 0$, so a \textbf{close-in reference distance} $d_0$ is used:
        \[ P_r(d) = P_r(d_0)\left(\frac{d_0}{d}\right)^2, \qquad d \ge d_0 \ge d_f \]
  \item $d_0$ is typically \textbf{1 km} for large cells, \textbf{100 m} for microcells,
        \textbf{1 m} indoors.
\end{itemize}}

\hr

\T{3.2 The three propagation mechanisms}

\Q{Explain the three basic radio wave propagation mechanisms. Large scale Vs small scale model}{\m{3+6}\m{3}\m{8}}
{\color{sub}\footnotesize \tF{3} \yr{\textbf{72 Ash}, \textbf{73 Bh}, 72 Ma}}

\sbsr{0.52}{%
\begin{itemize}
  \item Apart from the direct wave, the receiver gets \textbf{reflected}, \textbf{diffracted} and
        \textbf{scattered} waves.
  \item They add \textbf{vectorially}, so the resultant varies in strength in real time. That is
        the physical origin of fading.
\end{itemize}

\lead{1. Reflection}
\begin{itemize}
  \item Wave hits an object \textbf{very large compared to the wavelength}.
  \item Eg the earth's surface, buildings, walls.
  \item Part reflects, part transmits into the object.
  \item Strength set by the \textbf{reflection coefficient} $\Gamma$, which depends on
        polarisation, angle of incidence and the material.
  \item At grazing incidence $\Gamma \rightarrow -1$, ie \mk{full reflection with a
        $180^\circ$ phase flip}. That is the assumption the two-ray model rests on.
\end{itemize}}{%
\figT{c3_knife_edge.png}
\penalty0\vspace{1pt}
{\color{sub}\footnotesize Knife-edge diffraction geometry. Effective obstruction height $h$,
distances $d_1$ and $d_2$ from Tx and Rx.}}

\sbs{%
\lead{2. Diffraction \mk{(79 Bh pairs this with a Fresnel derivation, \m{2+6})}}
\begin{itemize}
  \item Path is obstructed by a surface with \textbf{sharp edges}.
  \item Secondary waves bend around the obstacle into the shadow region.
  \item \mk{This is what lets a mobile work with no LOS at all} in urban and hilly terrain.
  \item Explained by \textbf{Huygens' principle}: every point on a wavefront acts as a point
        source of secondary wavelets, which add to form the new front.
  \item Diffraction loss grows as the obstruction fills more of the first Fresnel zone.
\end{itemize}}{%
\lead{3. Scattering}
\begin{itemize}
  \item Medium contains objects \textbf{small compared to the wavelength}, or rough surfaces.
  \item Eg foliage, street signs, lamp posts, rain drops, snow.
  \item Energy is \textbf{spread in many directions}.
  \item Follows the same principles as diffraction.
  \item \mk{Received signal is often stronger than reflection and diffraction alone predict},
        because scattering fills in the gaps.
\end{itemize}}

\penalty0\vspace{3pt}
\sbsr{0.55}{%
\lead{Fresnel zones and the diffraction parameter}
\begin{itemize}
  \item A \textbf{Fresnel zone} is the ellipsoid of revolution over which secondary waves reach the
        receiver with an excess path of $n\lambda/2$.
  \item Numbered $F_1, F_2, F_3, \dots$. Odd zones add constructively, even zones destructively.
  \item Size set by \textbf{frequency} and \textbf{Tx--Rx distance}.
  \item \textbf{Fresnel--Kirchhoff diffraction parameter}
        \[ v \;=\; h\sqrt{\frac{2(d_1+d_2)}{\lambda d_1 d_2}} \]
  \item Radius of the $n$-th zone at distance $d_1$:
        \[ r_n = \sqrt{\frac{n\lambda d_1 d_2}{d_1+d_2}} \]
        \added
  \item \mk{Rule of thumb: keep 55\,\% of the first Fresnel zone clear} and the link behaves as
        free space.
  \item Diffraction gain $G_d$ (dB) against $v$:
  \begin{itemize}
    \item $v \le -1$: $G_d = 0$
    \item $-1 \le v \le 0$: $G_d = 20\log(0.5 - 0.62v)$
    \item $0 \le v \le 1$: $G_d = 20\log(0.5\,e^{-0.95v})$
    \item $v > 2.4$: $G_d = 20\log(0.225/v)$
  \end{itemize}
\end{itemize}}{%
\figT{c3_fresnel.png}
\penalty0\vspace{1pt}
{\color{sub}\footnotesize First and second Fresnel zones.}}

\hr

\T{3.3 Two-ray ground reflection model}

\Q{Derive the path difference and phase difference in the two-ray model, and the path loss}{\m{8}\m{2+6}\m{2+8}\m{6}}
{\color{sub}\footnotesize \tF{4} \yr{\textbf{71 Bh}, \textbf{74 Bh}, \textbf{76 Bh},
\textbf{\texttt{81 Bh}}} \gap$\cdot$\gap 76 Bh asks first for the \mk{(+) over free space}, so
open with that}

\sbsr{0.55}{%
\lead{Why it beats the free space model \mk{(the 2-mark opener)}}
\begin{itemize}
  \item Free space assumes \textbf{one} path. A real mobile link always has a
        \textbf{ground-reflected} path as well.
  \item The two-ray model adds that second path, so it predicts the
        \mk{$1/d^4$ roll-off actually measured} over large distances, instead of $1/d^2$.
  \item Found \textbf{reasonably accurate} for large-cell mobile systems over a few km.
  \item Needs no measured data, unlike Okumura or Hata.
\end{itemize}

\lead{Assumptions}
\begin{itemize}
  \item $h_t, h_r \gg \lambda$ and $h_t, h_r \ll d$.
  \item Reflected wave arrives at \textbf{grazing incidence}, so $\Gamma = -1$.
  \item Flat earth, no curvature.
\end{itemize}

\lead{Step 1: path difference $\Delta$ (method of images)}
\begin{itemize}
  \item Replace the reflection by an \textbf{image transmitter} at $-h_t$.
  \item LOS path and reflected path become two straight lines:
        \[ d' = \sqrt{(h_t-h_r)^2 + d^2}, \qquad d'' = \sqrt{(h_t+h_r)^2 + d^2} \]
  \item So
        \[ \Delta = d''-d' = d\sqrt{1+\left(\tfrac{h_t+h_r}{d}\right)^2}
                             - d\sqrt{1+\left(\tfrac{h_t-h_r}{d}\right)^2} \]
  \item Expand each root binomially, valid since $d \gg h_t,h_r$:
        \[ \Delta \approx d\left\{\left(1+\tfrac{1}{2}\left(\tfrac{h_t+h_r}{d}\right)^2\right)
           -\left(1+\tfrac{1}{2}\left(\tfrac{h_t-h_r}{d}\right)^2\right)\right\} \]
        \[ \boxed{\;\Delta \;\approx\; \frac{2 h_t h_r}{d}\;} \]
\end{itemize}

}{%
\figT{c3_tworay_geom.png}
\penalty0\vspace{2pt}
\figT{c3_tworay_images.png}
\penalty0\vspace{1pt}
{\color{sub}\footnotesize Method of images: the ground-reflected ray is redrawn as a straight
line from the image transmitter at $-h_t$.}}

\penalty0\vspace{3pt}
\sbs{%
\lead{Step 2: phase difference and time delay}
\[ \boxed{\;\theta_\Delta = \frac{2\pi\Delta}{\lambda} = \frac{4\pi h_t h_r}{\lambda d}\;} \]
\[ \tau_d = \frac{\Delta}{c} = \frac{\theta_\Delta}{2\pi f_c} \]

\lead{Step 3: total E-field}
\begin{itemize}
  \item For large $d$, $|E_0 d_0 / d'| \approx |E_0 d_0 / d''| \approx |E_0 d_0 / d|$, so the two
        rays differ \textbf{only in phase}:
        \[ |E_{\text{tot}}(d)| = \frac{E_0 d_0}{d}\left|\cos\theta_\Delta - \cos 0\right| \]
  \item Using $\cos\alpha - \cos\beta$ identity:
        \[ |E_{\text{tot}}(d)| = 2\,\frac{E_0 d_0}{d}\,\sin\!\left(\frac{\theta_\Delta}{2}\right) \]
  \item For $\theta_\Delta/2 < 0.3$ rad, $\sin x \approx x$:
        \[ \boxed{\;|E_{\text{tot}}(d)| \approx \frac{2 E_0 d_0}{d}\cdot
           \frac{2\pi h_t h_r}{\lambda d} = \frac{k}{d^2}\;} \]
  \item \mk{The field now falls as $1/d^2$, so power falls as $1/d^4$.}
\end{itemize}}{%
\lead{Step 4: received power and path loss}
\begin{itemize}
  \item Using $P_r = |E_{\text{tot}}|^2 A_e / \eta$:
        \[ \boxed{\;P_r \;=\; P_t G_t G_r \frac{h_t^2 h_r^2}{d^4}\;} \]
  \item \mk{Note $\lambda$ has cancelled}: the two-ray result is
        \textbf{independent of frequency}.
  \item In dB:
        \[ \begin{gathered} PL(\text{dB}) = 40\log d \\
           -\left(10\log G_t + 10\log G_r + 20\log h_t + 20\log h_r\right)
           \end{gathered} \]
  \item Valid only beyond the \textbf{crossover distance} $d_c = 4\pi h_t h_r/\lambda$;
        below it the model oscillates and free space is used instead. \added
\end{itemize}}

\hr

\T{3.4 Practical link budget design}

\Q{Explain the log-distance path loss model and log-normal shadowing}{\m{4}\m{2+4}}
{\color{sub}\footnotesize \tP{1} \yr{\textbf{77 Ch}} \gap$\cdot$\gap the \mk{link budget
numericals} built on it are 10 and 12 marks \yr{(\textbf{74 Bh}, \textbf{76 Bh})}
$\rightarrow$ Ch 3 companion}

\sbsr{0.53}{%
\lead{Log-distance path loss model}
\begin{itemize}
  \item Both theory and measurement show received power falls
        \textbf{logarithmically} with distance, indoors and outdoors.
  \item Generalises free space by letting the exponent be $n$ instead of 2:
        \[ \overline{PL}(d) \propto \left(\frac{d}{d_0}\right)^{n} \]
        \[ \boxed{\;\overline{PL}(d)\,[\text{dB}] = \overline{PL}(d_0)
           + 10\,n\log\!\left(\frac{d}{d_0}\right)\;} \]
  \item $\overline{PL}(d_0)$ is computed by the \textbf{free space} formula, or measured.
  \item Bars mean \textbf{ensemble average} over all possible paths at that distance.
  \item \mk{One number, $n$, absorbs the whole environment.}
\end{itemize}
\penalty0\vspace{2pt}
\figT{c3_pathloss_exp.png}}{%
\lead{Link budget design}
\begin{itemize}
  \item A link budget adds every gain and subtracts every loss between transmitter output and
        receiver input, then compares against the receiver \textbf{sensitivity}.
  \item Generic form, all terms in dB:
        \[ P_r = P_t + G_t + G_r - L_{\text{cable,t}} - L_{\text{cable,r}}
           - L_{\text{misc}} - PL(d) \]
  \item Link closes if $P_r \ge$ sensitivity $+$ \textbf{link margin}.
  \item \textbf{Two} budgets are needed, forward and reverse, and the
        \mk{weaker one sets the cell radius}.
\end{itemize}

\lead{How to run the exam version}
\begin{itemize}
  \item Compute \textbf{maximum allowable path loss} on each link:
        \[ PL_{\max} = P_t + G_t + G_r - L_{\text{cables}} - \text{sensitivity}
           - \text{margin} \]
  \item Put $PL_{\max}$ into the propagation model given (usually \textbf{Okumura}) and solve
        for $d$.
  \item Take the \textbf{smaller} of the forward and reverse distances.
  \item Compare against the range asked for and state whether the link is feasible.
\end{itemize}}

\penalty0\vspace{3pt}
\sbs{%
\lead{Log-normal shadowing}
\begin{itemize}
  \item Problem: two locations at the \textbf{same} $d$ can have wildly different surroundings,
        so a single average is wrong for any one site.
  \item Measurement shows $PL(d)$ at a given $d$ is \textbf{random}, and \textbf{log-normally}
        distributed about the average.
  \item Add a zero-mean Gaussian random variable $X_\sigma$ in dB, standard deviation $\sigma$ dB:
        \[ \boxed{\;PL(d) = \overline{PL}(d_0) + 10 n \log\!\left(\frac{d}{d_0}\right)
           + X_\sigma\;} \]
        \[ P_r(d)\,[\text{dBm}] = P_t\,[\text{dBm}] - PL(d) \]
  \item Called \textbf{shadowing}: the mobile is shadowed by clutter (buildings, hills, foliage).
  \item $n$ and $\sigma$ are found from measured data by \textbf{linear regression}, minimising
        mean square error.
  \item \mk{Small $\sigma$ $=$ the model fits that environment well.}
\end{itemize}}{%
\lead{Two ways to build a propagation model}
\begin{itemize}
  \item \textbf{Empirical}: collect measurements, fit curves. Eg Okumura, Hata.
  \begin{itemize}
    \item Accurate in the environment measured, no guarantee elsewhere.
  \end{itemize}
  \item \textbf{Analytical}: model the physics, derive equations. Eg free space, two-ray.
  \begin{itemize}
    \item Purely analytical models are inaccurate because so many assumptions are needed to keep
          the derivation tractable.
  \end{itemize}
  \item \mk{Practical models are hybrids}: an analytical core with measured correction factors.
\end{itemize}}

\hr

\T{3.5 Outdoor propagation models}

\Q{Explain any two outdoor propagation models. State the conditions for Okumura and Hata}{\m{3+3}\m{5}\m{2+2}\m{4}}
{\color{sub}\footnotesize \tS{7} \yr{71 Ma, 73 Ma, \textbf{75 Bh}, \textbf{79 Ch},
\textbf{80 Ch}, \textbf{\texttt{81 Bh}}, \textbf{\texttt{80 Bh}}} \gap$\cdot$\gap
\mk{the numerical version appears in 13 of 22 papers} $\rightarrow$ Ch 3 companion}

{\color{sub}\footnotesize Outdoor transmission runs over irregular terrain, so the
\textbf{terrain profile}, trees, buildings and hills must all be accounted for. Purely analytical
models cannot do this, so all the outdoor models below are \textbf{empirical}.}

\creamq{Okumura model}{\m{4}}
\sbsr{0.55}{%
\begin{itemize}
  \item \textbf{Wholly empirical.} Built from extensive drive-test measurements in and around
        Tokyo, 1968.
  \item \mk{The single most widely used model for signal prediction in urban areas.}
  \item Okumura published a set of \textbf{curves} of median attenuation relative to free space,
        $A_{mu}(f,d)$, for an urban area over quasi-smooth terrain, plus correction curves.
  \item \textbf{Validity conditions \mk{(the 2-mark answer)}}
  \begin{itemize}
    \item Frequency \textbf{150 MHz to 1920 MHz} (extrapolated to 3 GHz).
    \item Distance \textbf{1 km to 100 km}.
    \item BS antenna height \textbf{30 m to 1000 m}.
    \item Quasi-smooth terrain, urban reference.
  \end{itemize}
  \item \textbf{The formula}
        \[ \boxed{\;L_{50}(\text{dB}) = L_F + A_{mu}(f,d) - G(h_{te}) - G(h_{re})
           - G_{\text{AREA}}\;} \]
  \begin{itemize}
    \item $L_{50}$ $=$ median (50th percentile) path loss,
    \item $L_F$ $=$ free space path loss,
    \item $A_{mu}$ $=$ median attenuation relative to free space, read off the curve,
    \item $G(h_{te}), G(h_{re})$ $=$ antenna height gain factors,
    \item $G_{\text{AREA}}$ $=$ gain from the environment type.
  \end{itemize}
\end{itemize}}{%
\figT{c3_okumura_amu.png}
\penalty0\vspace{1pt}
{\color{sub}\footnotesize $A_{mu}(f,d)$ over a quasi-smooth urban area,
$h_{te} = 200$ m, $h_{re} = 3$ m.}}

\penalty0\vspace{3pt}
\sbsr{0.55}{%
\begin{itemize}
  \item \textbf{Height gain factors:}
        \[ G(h_{te}) = 20\log\frac{h_{te}}{200}, \qquad 30\,\text{m} < h_{te} < 1000\,\text{m} \]
        \[ G(h_{re}) = 10\log\frac{h_{re}}{3}, \qquad h_{re} \le 3\,\text{m} \]
        \[ G(h_{re}) = 20\log\frac{h_{re}}{3}, \qquad 3\,\text{m} < h_{re} < 10\,\text{m} \]
  \item \textbf{($+$)} simple, and \mk{the most accurate model available for mature cellular and
        land-mobile systems in cluttered urban areas}.
  \item \textbf{($-$)} slow response to rapid terrain change $\rightarrow$ good in cities, poorer
        in rural areas. Deviation from measurement is 10 to 14 dB.
  \item \textbf{($-$)} needs the curves, so it cannot be evaluated from formulas alone.
\end{itemize}}{%
\figT{c3_okumura_garea.png}
\penalty0\vspace{1pt}
{\color{sub}\footnotesize $G_{\text{AREA}}$ for suburban, quasi-open and open areas.}}

\creamq{Hata model}{\m{4}\m{6}\m{7}}
\sbsr{0.60}{%
\begin{itemize}
  \item \mk{Hata is Okumura's curves turned into formulas}, so it can be worked by hand.
  \item Same measured data, so the same physical assumptions.
  \item \textbf{Validity conditions}
  \begin{itemize}
    \item Frequency \textbf{150 MHz to 1500 MHz}.
    \item Distance \textbf{1 km to 20 km}.
    \item BS antenna \textbf{30 m to 200 m}, MS antenna \textbf{1 m to 10 m}.
    \item \textbf{No path-specific correction} is possible, unlike Okumura.
  \end{itemize}
  \item \textbf{Urban median path loss:}
  \begin{align*}
    L_{50}(\text{urban}) \;=\; & 69.55 + 26.16\log f_c - 13.82\log h_{te} \\
       & - a(h_{re}) + \left(44.9 - 6.55\log h_{te}\right)\log d
  \end{align*}
  with $f_c$ in MHz, $h_{te}, h_{re}$ in m, $d$ in km.
\end{itemize}}{%
\lead{Okumura Vs Hata \mk{(81 Bh asks exactly this, \m{4})}}
\begin{itemize}
  \item \textbf{Form}: Okumura is \textbf{graphical}, Hata is \textbf{closed-form}.
  \item \textbf{Frequency}: Okumura to 1920 MHz, Hata to 1500 MHz.
  \item \textbf{Distance}: Okumura 1 to 100 km, Hata 1 to 20 km.
  \item \textbf{BS height}: Okumura 30 to 1000 m, Hata 30 to 200 m.
  \item \textbf{Path-specific corrections}: Okumura \checkmark, Hata \xmark
  \item \textbf{Accuracy}: within \mk{1 dB of Okumura for $d > 1$ km}, so Hata is preferred for
        large-cell work. Okumura is better for small cells and irregular terrain.
  \item \textbf{Extension}: COST-231 extends Hata to \textbf{2 GHz} by changing the constants,
        for PCS bands. \added
\end{itemize}}

\penalty0\vspace{3pt}
\sbs{%
\begin{itemize}
  \item \textbf{Mobile antenna correction $a(h_{re})$:}
  \begin{itemize}
    \item small / medium city:
          $\;(1.1\log f_c - 0.7)h_{re} - (1.56\log f_c - 0.8)$ dB
    \item large city, $f_c \le 300$ MHz: $\;8.29(\log 1.54 h_{re})^2 - 1.1$ dB
    \item large city, $f_c > 300$ MHz: $\;3.2(\log 11.75 h_{re})^2 - 4.97$ dB
  \end{itemize}
  \item \textbf{Suburban:} $\;L_{50}(\text{urban}) - 2\left[\log(f_c/28)\right]^2 - 5.4$
  \item \textbf{Open rural:}
        $\;L_{50}(\text{urban}) - 4.78(\log f_c)^2 + 18.33\log f_c - 40.98$
  \item Suburban and rural terms are just \textbf{reductions in fixed loss} for less cluttered
        environments.
\end{itemize}}{%
\lead{Longley--Rice (ITS irregular terrain model)}
\begin{itemize}
  \item Point-to-point, \textbf{40 MHz to 100 GHz}, over irregular terrain.
  \item Runs as a \textbf{computer program}. Inputs: frequency, path length, polarisation,
        antenna heights, surface refractivity, ground conductivity and permittivity,
        climate.
  \item Two modes:
  \begin{itemize}
    \item \textbf{point-to-point}, using a detailed terrain profile,
    \item \textbf{area}, using empirical estimates where no profile exists.
  \end{itemize}
  \item \mk{($-$) ignores multipath, buildings and foliage}, so it is a terrain model, not an
        urban model.
\end{itemize}}

\penalty0\vspace{3pt}
\sbs{%
\lead{Walfisch--Bertoni \pill{gdL}{gd}{syllabus, no PYQ yet}}
\begin{itemize}
  \item Models \textbf{rooftop diffraction} explicitly, so it fits dense urban areas where the
        BS antenna is above the rooftops.
  \item Path loss is the product of three factors:
        \[ S = P_0 \, Q^2 \, P_1 \]
  \begin{itemize}
    \item $P_0$ $=$ free space loss between isotropic antennas,
    \item $Q^2$ $=$ reduction from the \textbf{row of buildings} shadowing the street,
    \item $P_1$ $=$ loss from the \textbf{final rooftop-to-street} diffraction.
  \end{itemize}
  \item Predicts an overall $\;\propto d^{-3.8}$ roll-off, matching urban measurement.
  \item Adopted as part of \textbf{COST-231}.
\end{itemize}}{%
\lead{Ericsson multiple breakpoint model \mk{(80 Bh asks by name, \m{2})}}
\begin{itemize}
  \item Built from measurements in a multi-floor office building.
  \item Uses a \textbf{uniform distribution} to generate path loss between a minimum and maximum
        bound at each distance.
  \item \textbf{4 breakpoints}, each carrying an upper and lower bound on path loss.
  \item Assumes \textbf{30 dB attenuation at $d_0 = 1$ m}, accurate at 900 MHz w/ unity-gain
        antennas.
  \item Slope changes from \textbf{20 dB/decade} at the first breakpoint to
        \textbf{60 dB/decade} in the final section.
  \item \mk{Its value is that it gives a deterministic limit on the range of path loss}, not a
        single number.
\end{itemize}}

\hr

\T{3.6 Indoor propagation models}

\Q{Explain indoor propagation models (any two). What factors influence them?}{\m{8}\m{5}}
{\color{sub}\footnotesize \tP{2} \yr{\textbf{70 Bh}, \textbf{\texttt{82 Bh}}}}

\sbsr{0.53}{%
\lead{How indoor differs from outdoor \mk{(the ``factors'' answer)}}
\begin{itemize}
  \item Distances covered are \textbf{much smaller}.
  \item Environment varies \textbf{far more} over a much smaller Tx--Rx separation.
  \item Propagation is governed by the \textbf{same three mechanisms}, but conditions are far more
        variable.
  \item \textbf{Factors that influence it:}
  \begin{itemize}
    \item \textbf{Layout} of the building, room sizes, corridors.
    \item \textbf{Construction materials}: concrete, metal, plaster, glass.
    \item \textbf{Building type}: office, residential, factory.
    \item Whether doors are \textbf{open or closed}.
    \item Where the antenna is mounted: \textbf{desk height Vs ceiling}.
    \item Number of \textbf{floors} between Tx and Rx.
    \item Number of \textbf{windows} and whether they are tinted (metallised tint is a strong
          reflector).
    \item Movement of \textbf{people and furniture}.
  \end{itemize}
  \item \mk{Path loss exponents indoors range from 1.6 (LOS corridor, better than free space,
        because of waveguiding) to 6 (obstructed).}
\end{itemize}}{%
\lead{Model 1: log-distance path loss (same form as outdoor)}
\[ PL(d)\,[\text{dB}] = \overline{PL}(d_0) + 10n\log\!\left(\frac{d}{d_0}\right) + X_\sigma \]
\begin{itemize}
  \item $n$ and $\sigma$ are tabulated \textbf{per building type}.
  \item Smaller $\sigma$ $\rightarrow$ better accuracy.
  \item $d_0$ is typically \textbf{1 m} indoors.
\end{itemize}

\lead{Model 2: attenuation factor model}
\[ \begin{gathered} PL(d)\,[\text{dB}] = \overline{PL}(d_0)
   + 10\,n_{SF}\log\!\left(\frac{d}{d_0}\right) \\
   + \text{FAF} + \sum \text{PAF} \end{gathered} \]
\begin{itemize}
  \item $n_{SF}$ $=$ path loss exponent measured on the \textbf{same floor}.
  \item \textbf{FAF} $=$ floor attenuation factor, added once per floor penetrated.
  \item \textbf{PAF} $=$ partition attenuation factor, summed over the partitions crossed.
  \item \mk{Standard deviation drops to about 4 dB}, against 13 dB for the plain
        log-distance model.
\end{itemize}}

\penalty0\vspace{3pt}
\sbs{%
\lead{Partition losses}
\begin{itemize}
  \item \textbf{Hard partitions}: structural walls of a room. Fixed.
  \item \textbf{Soft partitions}: movable dividers that do not reach the ceiling.
  \item \textbf{Between floors}: FAF is \mk{not linear in floor count}. The first floor costs the
        most (about 15 dB), each further floor adds only 6 to 10 dB, then it saturates.
  \item Reason: signal also travels \textbf{out of and back into} the building via windows,
        so the direct path stops dominating.
\end{itemize}}{%
\figT{c3_partition.png}
\penalty0\vspace{1pt}
{\color{sub}\footnotesize Measured same-floor partition losses.}}

\penalty0\vspace{3pt}
\sbs{%
\lead{Model 3: ITU indoor path loss model \added}
\[ L_{total}\,[\text{dB}] = 20\log(f) + N\log(d) + L_f(n) - 28 \]
\begin{itemize}
  \item $f$ $=$ frequency in \textbf{MHz}; $d$ $=$ separation in \textbf{metres}, $d > 1$ m.
  \item $N$ $=$ \textbf{distance power loss coefficient} (tabulated per band and building type;
        note it plays the role of $10n$, so $N \approx 30$ for an office at 1.8 GHz).
  \item $L_f(n)$ $=$ \textbf{floor penetration loss factor} for $n$ floors crossed, read from a
        table rather than computed.
  \item \mk{Why it looks different}: the $20\log f$ term and the $-28$ constant are the free
        space reference folded in, so no separate $\overline{PL}(d_0)$ measurement is needed.
        \mk{That is its selling point} $-$ it is fully \textbf{empirical / tabulated}, usable
        without a site survey.
\end{itemize}}{%
\lead{Which two to write \mk{(the question always says ``any two'')}}
\begin{itemize}
  \item Safest pair: \textbf{log-distance} $+$ \textbf{attenuation factor}. They share a form,
        so one derivation covers both, and FAF/PAF is where the marks are.
  \item \textbf{ITU} is the cheapest third model to add: one formula, three symbols, no
        derivation.
  \item \mk{The Ericsson multiple breakpoint model is also an indoor model} $-$ it is written up
        in \S3.5 with the outdoor set because \textbf{80 Bh} asks for it by name there. Use it
        here if the paper wants a model that gives \textbf{bounds} rather than a mean. \added
  \item \mk{Never write only the formula.} Each model scores on
        \textbf{what it adds over free space}: $n \ne 2$, then FAF/PAF, then floor tables.
\end{itemize}}

\hr

\T{3.7 Small-scale fading and multipath}

\Q{What is small-scale fading? Give its types and the factors that influence it}{\m{2+4+4}\m{4+4}\m{3+6}\m{1+2+5}\m{2+2}}
{\color{sub}\footnotesize \tS{6} \yr{72 Ma, 73 Ma, \textbf{75 Bh}, \textbf{78 Ch},
\texttt{80 Ba}, \texttt{81 Ba}} \gap$\cdot$\gap
\mk{the most repeated question in the chapter}}

\sbsr{0.53}{%
\lead{Definition}
\begin{itemize}
  \item \textbf{Small-scale fading} $=$ rapid fluctuation of the amplitude, phase and multipath
        delay of a radio signal over a \textbf{short travel distance or short time interval}.
  \item Caused by \textbf{interference between two or more copies} of the transmitted signal that
        arrive at slightly different times.
  \item The copies (multipath waves) add \textbf{vectorially} with random phases, so the sum
        behaves like noise.
  \item \mk{Received power can change by 30 to 40 dB when the receiver moves a fraction of a
        wavelength.}
  \item Large-scale path loss can be ignored over such a short distance.
\end{itemize}

\lead{The three effects it produces}
\begin{itemize}
  \item \textbf{Rapid changes in signal strength} over a small travel distance or time interval.
  \item \textbf{Random frequency modulation} due to different \textbf{Doppler shifts} on
        different multipath components.
  \item \textbf{Time dispersion (echoes)} caused by multipath propagation delays
        $\rightarrow$ \mk{ISI}.
\end{itemize}}{%
\lead{Time dispersion Vs frequency dispersion}
\begin{itemize}
  \item \textbf{Time dispersion}: received signal \textbf{lasts longer} than the transmitted one,
        because paths have different delays. Measured by \textbf{delay spread}.
  \item \textbf{Frequency dispersion}: received signal is \textbf{wider in frequency} than the
        transmitted one, because paths have different Doppler shifts. Measured by
        \textbf{Doppler spread}.
  \item Both distort the received signal, and they are \textbf{independent} of each other.
        \mk{That is why fading is classified twice over, once by each.}
\end{itemize}

\lead{Impulse response model of the multipath channel \pill{gdL}{gd}{syllabus, no PYQ yet}}
\begin{itemize}
  \item The channel is treated as a \textbf{linear filter with a time-varying impulse response}.
  \item Feed it a unit impulse $\delta(t)$ and the output is
        \[ h(t) \;=\; \sum_{n=1}^{N} A_n\, e^{-j\phi_n}\,\delta(t-\tau_n) \]
  \item $A_n$, $\tau_n$, $\phi_n$ $=$ amplitude, arrival delay and phase of the $n$-th path.
  \item Time variation comes from motion, so it is really $h(t,\tau)$: $\tau$ indexes the
        multipath delay, $t$ indexes the mobile's position.
  \item \mk{This one model generates everything else in the chapter}: the PDP is
        $|h(t,\tau)|^2$, and the delay-spread parameters are its moments.
\end{itemize}}

\penalty0\vspace{3pt}
\sbs{%
\lead{Factors influencing small-scale fading \mk{(learn all 4)}}
\begin{itemize}
  \item \textbf{Multipath propagation.}
  \begin{itemize}
    \item Random amplitude and phase of each component $\rightarrow$ the sum fluctuates.
    \item Also \textbf{lengthens} the time the baseband signal takes to arrive
          $\rightarrow$ signal smearing $\rightarrow$ ISI.
  \end{itemize}
  \item \textbf{Speed of the mobile.}
  \begin{itemize}
    \item Relative motion causes a \textbf{Doppler shift} $f_d$ on each component.
    \item Different components arrive from different angles $\rightarrow$ different shifts
          $\rightarrow$ random FM.
  \end{itemize}
  \item \textbf{Speed of surrounding objects.}
  \begin{itemize}
    \item If nearby objects move \textbf{faster than the mobile}, they dominate the fading.
    \item If they move slower, ignore them.
  \end{itemize}
  \item \textbf{Transmission bandwidth of the signal.}
  \begin{itemize}
    \item If signal bandwidth $>$ channel bandwidth, the received signal is
          \textbf{distorted}, but its strength does not fade much over a local area.
    \item If signal bandwidth $<$ channel bandwidth, amplitude \textbf{changes rapidly}
          but the signal is not distorted in time.
    \item \mk{A channel gives you distortion or fading, typically one or the other.}
  \end{itemize}
\end{itemize}}{%
\lead{Multipath measurement techniques \pill{gdL}{gd}{syllabus, no PYQ yet}}
\begin{itemize}
  \item \textbf{Direct RF pulse sounding}: transmit a narrow pulse, display the received power
        against delay on a scope. Simple, but a single pulse can be lost in a fade.
  \item \textbf{Spread spectrum sliding correlator}: transmit a PN sequence, correlate at the
        receiver w/ a slightly slower clock so the correlation peak slides through the delays.
        Good dynamic range, rejects interference.
  \item \textbf{Swept frequency (frequency domain) sounding}: sweep a network analyser across the
        band, measure $S_{21}$, then inverse Fourier transform to get $h(\tau)$.
        Needs a cabled link, so it suits indoor work only.
\end{itemize}}

\Q{Define Doppler shift and Doppler spread. Derive the expression for Doppler shift}{\m{2+4}\m{1+2+5}\m{2}}
{\color{sub}\footnotesize \tF{4} \yr{\textbf{78 Ch}, \texttt{81 Ba}, \textbf{80 Ch},
\textbf{\texttt{82 Bh}}} \gap$\cdot$\gap 82 Bh words it as
``significance of frequency dispersion''}

\sbsr{0.55}{%
\lead{Derivation \mk{(reproduce these 5 lines exactly)}}
\begin{itemize}
  \item Mobile moves at constant velocity $v$ along a segment of length $d$ from $X$ to $Y$,
        receiving from a \textbf{distant} source $S$.
  \item Source is far away $\rightarrow$ the arrival angle $\theta$ is the same at $X$ and $Y$.
  \item \textbf{Path length difference}
        \[ \Delta l \;=\; d\cos\theta \;=\; v\,\Delta t\,\cos\theta \]
  \item \textbf{Phase change} in the received signal
        \[ \Delta\phi \;=\; \frac{2\pi \Delta l}{\lambda}
           \;=\; \frac{2\pi v \Delta t}{\lambda}\cos\theta \]
  \item \textbf{Apparent frequency change} is the rate of change of phase:
        \[ f_d \;=\; \frac{1}{2\pi}\cdot\frac{\Delta\phi}{\Delta t}
           \;\Longrightarrow\;
           \boxed{\;f_d \;=\; \frac{v}{\lambda}\cos\theta\;} \]
\end{itemize}}{%
\figT{c3_doppler.png}
\penalty0\vspace{1pt}
{\color{sub}\footnotesize Mobile moving from $X$ to $Y$ at velocity $v$; source $S$ is far away,
so $\theta$ is unchanged.}}

\penalty0\vspace{3pt}
\sbs{%
\lead{Reading the result}
\begin{itemize}
  \item $\theta = 0^\circ$, moving \textbf{towards} the source: $f_d = +v/\lambda$, maximum
        \textbf{positive} shift. Received frequency $f_c + f_d$.
  \item $\theta = 180^\circ$, moving \textbf{away}: $f_d = -v/\lambda$, maximum negative shift.
  \item $\theta = 90^\circ$, moving \textbf{perpendicular}: \mk{$f_d = 0$}, no shift at all.
  \item \textbf{Maximum Doppler shift}
        \[ f_m = \frac{v}{\lambda} = \frac{v f_c}{c} \]
\end{itemize}

\lead{Doppler spread $B_D$}
\begin{itemize}
  \item Measure of the \textbf{spectral broadening} caused by the time rate of change of the
        channel.
  \item Defined as the \textbf{maximum Doppler shift}: $B_D = f_m = v/\lambda$.
  \item A pure tone at $f_c$ comes out spread over $f_c \pm f_m$.
  \item Characterises the channel's \textbf{frequency dispersiveness}.
\end{itemize}}{%
\lead{Significance of frequency dispersion \mk{(82 Bh)}}
\begin{itemize}
  \item It makes the channel \textbf{time-varying}, so the receiver must keep re-estimating it.
  \item If $B_D$ approaches the signal bandwidth, the channel changes \textbf{within one symbol}
        $\rightarrow$ \mk{fast fading, and an irreducible BER floor that more transmit power
        cannot cure}.
  \item It sets how often an \textbf{adaptive equalizer must retrain} and how fast a
        \textbf{carrier recovery loop} must track.
  \item Different multipath components have different $f_d$, so the composite signal suffers
        \textbf{random FM}, heard as noise on an analog system.
\end{itemize}}

\Q{Explain the parameters of the mobile multipath channel}{\m{7}\m{4}\m{2+4}}
{\color{sub}\footnotesize \tF{3} \yr{71 Ma, \textbf{\texttt{80 Bh}}, 70 Ma}
\gap$\cdot$\gap 80 Bh asks which parameters are used to \mk{classify the types of fading}
\gap$\cdot$\gap \mk{numericals on this appear in 5 papers} $\rightarrow$ Ch 3 companion}

\sbsr{0.55}{%
\lead{Power delay profile (PDP)}
\begin{itemize}
  \item \textbf{PDP} $=$ received power against \textbf{excess delay}, ie the intensity of the
        signal arriving through the multipath channel as a function of delay.
  \item Obtained by \textbf{averaging many impulse response measurements} over a local area.
  \item Every time-dispersion parameter below is a \textbf{moment of the PDP}.
\end{itemize}}{%
\figT{c3_pdp.png}
\penalty0\vspace{1pt}
{\color{sub}\footnotesize A measured PDP w/ its mean excess delay, rms delay spread and 10 dB
maximum excess delay marked. Note the $-20$ dB noise threshold.}}

\penalty0\vspace{3pt}
\sbs{%
\lead{Time dispersion parameters}
\begin{itemize}
  \item \textbf{First arrival delay $\tau_A$}: delay of the first signal to reach the receiver,
        set by the shortest path. Used as the \textbf{reference}, so every other delay is an
        \emph{excess} delay.
  \item \textbf{Mean excess delay} $=$ first moment of the PDP:
        \[ \bar{\tau} \;=\; \frac{\sum_k a_k^2 \tau_k}{\sum_k a_k^2}
                       \;=\; \frac{\sum_k P(\tau_k)\,\tau_k}{\sum_k P(\tau_k)} \]
  \item \textbf{rms delay spread} $=$ square root of the second central moment:
        \[ \sigma_\tau = \sqrt{\overline{\tau^2} - (\bar{\tau})^2},
           \qquad \overline{\tau^2} = \frac{\sum_k a_k^2 \tau_k^2}{\sum_k a_k^2} \]
  \item \textbf{Maximum excess delay (X dB)} $=$ $\tau_X - \tau_0$, the delay over which multipath
        energy stays within X dB of the strongest component.
  \item \mk{$\sigma_\tau$ is the one that matters in practice}, because it, not $\tau_{\max}$,
        predicts ISI.
  \item These values \textbf{depend on the noise threshold} chosen. Set the threshold too low and
        noise is counted as multipath, inflating every parameter.
\end{itemize}}{%
\lead{Coherence bandwidth $B_c$}
\begin{itemize}
  \item \textbf{Definition}: the range of frequencies over which the channel passes all spectral
        components with \textbf{approximately equal gain and linear phase}, ie over which the
        channel looks \textbf{flat}.
  \item It is the \textbf{frequency-domain dual} of the rms delay spread. There is no exact
        relationship, only rules of thumb:
        \[ B_c \approx \frac{1}{50\,\sigma_\tau} \quad (\text{0.9 correlation, ``90\,\%''}) \]
        \[ B_c \approx \frac{1}{5\,\sigma_\tau} \quad (\text{0.5 correlation, ``50\,\%''}) \]
  \item Two sinusoids separated by \textbf{more} than $B_c$ are affected \textbf{differently} by
        the channel.
  \item \textbf{Significance:} compare $B_c$ against the signal bandwidth to decide whether the
        channel is \mk{flat or frequency selective}. That single comparison decides whether an
        equalizer is needed.
\end{itemize}}

\penalty0\vspace{2pt}
\lead{Coherence time $T_c$}
\begin{itemize}
  \item \textbf{Definition}: the time over which the channel impulse response is essentially
        \textbf{time-invariant}.
  \item Dual of the Doppler spread:
        \[ T_c \approx \frac{1}{f_m}, \qquad
           T_c \approx \sqrt{\frac{9}{16\pi f_m^2}} = \frac{0.423}{f_m} \]
        The second (geometric mean) form is the one normally used.
  \item Two signals arriving more than $T_c$ apart are affected \textbf{differently}.
  \item \textbf{Large $T_c$} $\rightarrow$ channel changes slowly.
  \item \textbf{Significance:} compare $T_c$ against the symbol period to decide whether the
        channel is \mk{slow or fast fading}. Also sets how often the equalizer must retrain.
\end{itemize}

\Q{Explain the types of small-scale fading}{\m{4+4}\m{1+3}\m{2+7}\m{2+4}}
{\color{sub}\footnotesize \tS{6} \yr{70 Ma, 74 Ma, \textbf{79 Ch}, \texttt{81 Ba},
\textbf{80 Ch}, \textbf{\texttt{80 Bh}}} \gap$\cdot$\gap
\mk{fading is classified twice over: once by delay spread, once by Doppler spread. Always give
both tables.}}

\sbs{%
\lead{A. By multipath time delay spread}
\penalty0\vspace{2pt}
\begin{tabularx}{\linewidth}{@{}L{2.05cm}XX@{}}
\toprule
\hrow & \thead{Flat fading} & \thead{Frequency selective} \\
\midrule
Bandwidth & $B_S \ll B_C$ & $B_S > B_C$ \\
Delay spread & $\sigma_\tau \ll T_S$ & $\sigma_\tau > T_S$ \\
Spectrum of the signal & \textbf{preserved} & \textbf{not preserved} \\
Channel acts as & flat gain, linear phase & a filter \\
ISI & none & \mk{severe} \\
Cure & raise transmit power, use diversity & \mk{equalizer} \\
\bottomrule
\end{tabularx}
\penalty0\vspace{3pt}
\begin{itemize}
  \item \textbf{Flat fading} is the most common type.
  \item Signal bandwidth is narrow, so the whole signal rides up and down together.
  \item Causes \mk{deep fades}, so flat-fading channels are sometimes called
        \textbf{amplitude varying} or \textbf{narrowband} channels.
  \item \textbf{Frequency selective fading} occurs when the delayed copies overlap
        \textbf{into the next symbol}.
  \item \mk{Frequency selective fading is harder to model}, because the channel response depends
        on the exact multipath structure.
\end{itemize}}{%
\lead{B. By Doppler spread}
\penalty0\vspace{2pt}
\begin{tabularx}{\linewidth}{@{}L{2.05cm}XX@{}}
\toprule
\hrow & \thead{Fast fading} & \thead{Slow fading} \\
\midrule
Doppler spread & high, $B_S < B_D$ & low, $B_S \gg B_D$ \\
Coherence time & $T_C < T_S$ & $T_C \gg T_S$ \\
Channel changes & \textbf{within} a symbol & over \textbf{many} symbols \\
Relative rates & channel varies \textbf{faster} than the baseband signal & channel varies
  \textbf{slower} than the baseband signal \\
\bottomrule
\end{tabularx}
\penalty0\vspace{3pt}
\begin{itemize}
  \item Fast / slow says \textbf{nothing} about flat / selective. \mk{The four combinations are all
        possible}: flat slow, flat fast, frequency selective slow, frequency selective fast.
  \item \textbf{Fast fading} arises from motion, or from motion of the surroundings, not from
        distance.
  \item A fast fading channel produces an \mk{irreducible BER floor}: extra transmit power does
        not help, because the channel has already changed by the time the symbol ends.
  \item In practice fast fading only matters for \textbf{very low data rates}.
\end{itemize}}

\hr

\T{3.8 Rayleigh and Ricean fading}

\Q{Rayleigh fading channel Vs Ricean fading channel, with appropriate expressions}{\m{5}\m{2}\m{2+4}}
{\color{sub}\footnotesize \tF{5} \yr{70 Ma, 71 Ma, \textbf{72 Ash}, \textbf{73 Bh},
\textbf{80 Ch}}}

\lead{What these distributions describe}
\begin{itemize}
  \item Not \emph{when} the signal fades, but \textbf{how the received envelope is statistically
        distributed}.
  \item Both apply to \textbf{flat fading} channels.
\end{itemize}

\sbs{%
\lead{Rayleigh}
\begin{itemize}
  \item Applies when \textbf{all multipath components have roughly equal amplitude} and there is
        \mk{no dominant LOS component}.
  \item Physically: mobile deep in an urban area, far from the BS, no line of sight.
  \item The in-phase and quadrature sums are each Gaussian, so the \textbf{envelope} is Rayleigh:
        \[ p(r) = \begin{cases}
           \dfrac{r}{\sigma^2}\,e^{-r^2/2\sigma^2} & 0 \le r \le \infty \\[4pt]
           0 & r < 0 \end{cases} \]
  \item $\sigma$ $=$ rms value of the received voltage \textbf{before} envelope detection.
  \item $\sigma^2$ $=$ time-average power before envelope detection.
  \item \textbf{Phase} is uniformly distributed over $0$ to $2\pi$.
  \item Key values:
  \begin{itemize}
    \item mean $\;r_{\text{mean}} = \sigma\sqrt{\pi/2} = 1.2533\,\sigma$
    \item median $\;r_{\text{median}} = 1.177\,\sigma$
    \item variance $\;\sigma_r^2 = 0.4292\,\sigma^2$
  \end{itemize}
  \item \mk{Mean and median differ by only 0.55 dB}, which is why the median is used in practice.
\end{itemize}}{%
\lead{Ricean}
\begin{itemize}
  \item Applies when there \textbf{is} a \textbf{stationary, non-fading, dominant} component,
        normally a LOS path, on top of the random multipath.
  \item Physically: microcell, indoor pico-cell, or a mobile with clear sight of the BS.
  \item The random components add onto a fixed phasor, so:
        \[ p(r) = \begin{cases}
           \dfrac{r}{\sigma^2}\,e^{-(r^2+A^2)/2\sigma^2}\,
             I_0\!\left(\dfrac{Ar}{\sigma^2}\right) & A \ge 0,\; r \ge 0 \\[4pt]
           0 & r < 0 \end{cases} \]
  \item $A$ $=$ peak amplitude of the \textbf{dominant} signal.
  \item $I_0(\cdot)$ $=$ modified Bessel function of the first kind, zeroth order.
  \item \textbf{Ricean $K$-factor}
        \[ K = \frac{A^2}{2\sigma^2}, \qquad K\,[\text{dB}] = 10\log\frac{A^2}{2\sigma^2} \]
        $=$ ratio of \textbf{deterministic power to multipath power}.
  \item \mk{$A \rightarrow 0$, ie $K \rightarrow -\infty$ dB: Ricean degenerates into Rayleigh.}
  \item $K \gg 1$: the distribution tends to \textbf{Gaussian} about the mean, ie almost no
        fading.
\end{itemize}}

\sbs{%
\figT{c3_rayleigh_pdf.png}
\penalty0\vspace{1pt}
{\color{sub}\footnotesize Rayleigh envelope pdf.}}{%
\figT{c3_ricean_pdf.png}
\penalty0\vspace{1pt}
{\color{sub}\footnotesize $K = -\infty$ dB is the Rayleigh case; $K = 6$ dB shows a dominant
component pulling the distribution towards Gaussian.}}
